\subsection{Работа с модулями платформы}
- \verb|Ledbar.new(N)| создаёт объект управления светодиодами (используем \verb|N=25|).
- \verb|leds:set(i, r, g, b)| задаёт цвет по индексу \verb|i| (\verb|r,g,b| в \verb|[0.0,1.0]|).
- \verb|Timer.new(period, fn)| создаёт периодический таймер; \verb|Timer.callLater(delay, fn)| — отложенный вызов.
- \verb|function callback(event)| — обработчик системных событий.

\subsection{Типичные ошибки новичка}
\begin{enumerate}
    \item \textbf{Регистр букв}: \texttt{Print}, \texttt{PRINT} и \texttt{print} — это три разных слова. В Lua все ключевые слова и функции пишутся строчными буквами (кроме классов вроде \texttt{Ledbar}).
    \item \textbf{Путаница с индексами}:
    \begin{itemize}
        \item Обычные таблицы Lua нумеруются с \textbf{1}: \texttt{t[1]}.
        \item Аппаратные модули (светодиоды, массивы из C++) часто нумеруются с \textbf{0}: \texttt{leds:set(0, ...)}. Будьте внимательны!
    \end{itemize}
    \item \textbf{Бесконечные циклы}:
    \begin{minted}{lua}
-- ОПАСНО! Этот код "повесит" выполнение скрипта
while true do
    print("Я завис")
end
    \end{minted}
    Вместо вечных циклов используйте \texttt{Timer}.
    \item \textbf{Ошибка "attempt to index a nil value"}: Вы пытаетесь обратиться к полю переменной, которая пуста (\texttt{nil}). Проверьте, правильно ли вы создали объект или таблицу.
\end{enumerate}

\subsection{Примеры, используемые далее}
\begin{infobox}[Обратите внимание]
    В следующих примерах используется библиотека \texttt{Ledbar}, которую мы подробно изучим в следующей главе. Сейчас сосредоточьтесь на синтаксисе языка Lua (циклы, функции, таблицы), не вдаваясь в детали работы светодиодов.
\end{infobox}
\begin{minted}{lua}
local unpack = table.unpack
local ledNumber = 25
local leds = Ledbar.new(ledNumber)
local function changeColor(col)
  for i=0, ledNumber-1 do
    leds:set(i, unpack(col))
  end
end
changeColor({0,1,0})
\end{minted}
\begin{minted}{lua}
local ledNumber = 25
local leds = Ledbar.new(ledNumber)
for i=0, ledNumber-1 do leds:set(i, 0, 0, 0) end
Timer.callLater(3, function ()
  for i=0, ledNumber-1 do leds:set(i, 0, 0, 1) end
end)
\end{minted}
\begin{minted}{lua}
local ledNumber = 25
local leds = Ledbar.new(ledNumber)
on = false
function updateBlink()
  local r = on and 1 or 0
  for i=0, ledNumber-1 do leds:set(i, r, 0, 0) end
  on = not on
end
blinkTimer = Timer.new(0.5, updateBlink)
blinkTimer:start()
\end{minted}
